\section{Introduction}\label{sec:intro}

In a series of papers motivated by molecular spaces and digital topology, Ivashchenko~\cite{Ivashchenko1993,Ivashchenko1994a,Ivashchenko1994b} introduced four \emph{contractible transformations} of a finite simple graph: deleting or adding a vertex whose open neighbourhood is itself contractible, and deleting or adding an edge $uv$ whose common neighbourhood $\Nb(u)\cap\Nb(v)$ is contractible. The graphs that can be reduced to a single vertex by such transformations form the class $\I$ of \emph{I-contractible} graphs. Chen, Yau and Yeh~\cite{ChenYauYeh2001} proved that contractible transformations preserve the simple homotopy type of the clique complex $\cl(G)$, so every graph in $\I$ has a contractible clique complex; very recently Dochtermann and Matsushita~\cite{DochtermannMatsushita2025} proved the converse, so that $G\in\I$ if and only if $\cl(G)$ is contractible, and that I-homotopy of graphs coincides with simple homotopy equivalence of clique complexes. Contractible transformations have since been connected to strong homotopy types~\cite{BarmakMinian2012}, to simplicial simple-homotopy of flag complexes~\cite{BouletFieuxJouve2010,BouletFieuxJouve2020corr}, to dismantlability hierarchies~\cite{FieuxJouve2020}, and to collapsibility questions for flag complexes~\cite{EFH2022,Malen2023}.

Because the definition of $\I$ allows edges and vertices to be \emph{added}, deciding membership is delicate. Two natural subclasses are easier to handle~\cite{DEFH2023}: the \emph{strong I-contractible} graphs $\IS$, which can be reduced to a vertex using deletions only (of vertices with $\Nb(v)\in\IS$ and edges with $\Nb(u,v)\in\IS$), and the \emph{vertex-strong I-contractible} graphs $\IVS$, which can be reduced using vertex deletions only. One has
\[
  \IVS\ \subseteq\ \IS\ \subseteq\ \I\ \subseteq\ \{G : \cl(G)\ \text{is contractible}\},
\]
with the last inclusion an equality by~\cite{DochtermannMatsushita2025}. This is the graph-theoretic shadow of the classical chain \emph{non-evasive} $\subset$ \emph{collapsible} $\subset$ \emph{contractible} for simplicial complexes~\cite{Jonsson2008,Cohen1973}, and, as for complexes, the interesting questions concern the gaps. Dochtermann, Espinoza, Fr\'ias-Armenta and Hern\'andez~\cite{DEFH2023} determined by computer the smallest members of $\IS\setminus\IVS$ (eleven vertices), computed the cardinalities of all these classes for connected graphs on at most eleven vertices, and exhibited a $15$-vertex flag triangulation $\Dunce$ of Zeeman's dunce hat~\cite{Zeeman1963} which is I-contractible but not strong I-contractible (Fig.~2A of~\cite{DEFH2023}). Since $\IS$ is contained in the class of graphs whose clique complex is collapsible~\cite{EFH2022}, and the dunce hat is the standard example of a contractible non-collapsible complex, this is a natural candidate for a smallest example, and the authors conjectured:

\begin{quote}
\textbf{Conjecture 14 of~\cite{DEFH2023}.} \emph{The smallest graph (in terms of vertices) with trivial $\Ztwo$-homology that is not strong I-contractible is shown in Fig.~2A. In particular this is the smallest graph that is I-contractible but not strong I-contractible.}
\end{quote}

A companion statement, Conjecture~15 of~\cite{DEFH2023}, asserts that the graph of Fig.~2B (the dunce hat with one extra edge) is the smallest strong I-contractible graph $G$ having an edge $\{v,w\}$ with $\Nb(v,w)\in\IS$ such that $G\setminus\{v,w\}$ is not strong I-contractible---that is, the smallest graph in which a ``legal'' strong deletion of an edge destroys strong I-contractibility.

In this paper we settle both conjectures in the negative by an exhaustive and independently certified computer search, and we find that the two sentences of Conjecture~14 fail for \emph{different} reasons.

\paragraph{Main results} Let $\cl(G)$ denote the clique complex. A graph is \emph{$\Ztwo$-acyclic} if $\widetilde H_*(\cl(G);\Ztwo)=0$.

\begin{itemize}[leftmargin=2em]
\item[(A)] (\emph{Thirteen vertices, torsion}; Theorem~\ref{thm:A}.) Every $\Ztwo$-acyclic graph on at most $12$ vertices is strong I-contractible. On $13$ vertices there are exactly three $\Ztwo$-acyclic graphs that are not strong I-contractible. Their clique complexes have $H_1(\,\cdot\,;\ZZ)\cong\ZZ/3$ and $H_2=0$: they are $2$-dimensional Moore spaces $M(\ZZ/3,1)$, hence not contractible and in particular not I-contractible. The first sentence of Conjecture~14 is therefore false, but not because of a smaller dunce hat: it is false because $\Ztwo$-homology does not see odd torsion. By Adamaszek's classification of torsion in flag complexes with at most $12$ vertices~\cite{Adamaszek2014}, these are the smallest flag complexes with odd torsion in $H_1$.
\item[(B)] (\emph{Fourteen vertices, the gap $\I\setminus\IS$}; Theorem~\ref{thm:B}.) No graph on at most $13$ vertices is I-contractible but not strong I-contractible, and there are such graphs on $14$ vertices. Among $K_4$-free graphs on $14$ vertices (equivalently, graphs whose clique complex is $2$-dimensional, the class containing the dunce hat) there are exactly $51$ \emph{stuck} $\Ztwo$-acyclic graphs that are not in $\IS$: three of them, $\Ga,\Gb,\Gc$, are contractible and I-contractible, and the remaining $48$ have $H_1\cong\ZZ/3$. The second sentence of Conjecture~14 is therefore false, with exact minimum $14$ instead of $15$.
\item[(C)] (\emph{Reduction}; Theorem~\ref{thm:C}.) A smallest $\Ztwo$-acyclic graph outside $\IS$ with the minimum number of edges is \emph{stuck}: no open neighbourhood $\Nb(v)$ and no non-empty common neighbourhood $\Nb(u,v)$ is $\Ztwo$-acyclic. Stuck graphs on at most $22$ vertices are $2$-connected, have minimum degree at least $4$ and no vertex of degree $n-1$. This reduces the search space by three to four orders of magnitude and is what makes the enumeration feasible.
\item[(D)] (\emph{Conjecture~15}; Theorem~\ref{thm:D}.) The smallest strong I-contractible graph with an edge $\{v,w\}$, $\Nb(v,w)\in\IS$, whose deletion leaves a graph outside $\IS$ has $14$ vertices: $\Ga+\{0,1\}$ is such a graph, and there is none on $13$ vertices.
\item[(E)] (\emph{Table}; Section~\ref{sec:table}.) We recompute the cardinality table of~\cite{DEFH2023} for $n\le 11$, obtaining identical values in all cells, add the column $\I$, and extend the table to $n=12$ (connected graphs: $\num{164059830476}$) with the new value $|\IVS(12)|=\num{22666525713}$; rows for $n=13$ and for $K_4$-free graphs on $14$ vertices are in progress and reported where available.
\end{itemize}

\paragraph{Methodological remark} The examples of (A) show that $\Ztwo$-homology, which is the natural invariant to compute in a search of this size, is not sufficient in this problem: the integer homology of candidates must be computed, and the finite certificates we provide (deletion sequences, Smith normal forms) are over $\ZZ$. Our examples in (B) are, like $\Dunce$, flag triangulations of a disc with its boundary identified along a short cycle with the dunce-hat pattern $a\,a\,a^{-1}$; the $48$ torsion examples realise the pattern $a^3$ on cycles of length $4$, $5$ and $6$. We conjecture that this picture persists (Section~\ref{sec:open}).

\paragraph{Reproducibility} All code, the exact cluster scripts, the raw logs, every graph found and human-readable certificates are available in a public repository,
\begin{center}\small\repo\end{center}
A script \texttt{verify.sh} re-checks every finite certificate behind Theorems~\ref{thm:A}, \ref{thm:B} and~\ref{thm:D} with a dependency-free Python implementation, independent of the C code used in the search, in under a minute on a laptop.

\paragraph{Organisation} Section~\ref{sec:prelim} recalls the definitions. Section~\ref{sec:reduction} proves the reduction to stuck graphs. Section~\ref{sec:algorithms} describes the algorithms, the distributed computation and its validation. Section~\ref{sec:table} presents the extended table. Section~\ref{sec:examples} describes the minimal examples in detail. Section~\ref{sec:consequences} discusses consequences (including the conjectures of~\cite{FriasArmenta2020}) and open problems. The appendices contain the proofs of the reduction lemmas, the list of the $51$ stuck graphs, the complete certificates and the cluster accounting.
